% ===== 原文 PDF pp.86--88（印刷 pp.85--87）=====
% tr-08 / 块 tag 5.28--5.36 / 文件 a
% 起点：p.86「We thus have a Class S description …」——p.85 上开始、跨到 p.86 前三行的那半句
%       （"…rise to the type G singularity … fibers over the y-plane."）由 tr-07 按「开始者写完」收尾，本文件不重译。
% 终点：p.88 末（含一句跨到 p.89 的句子，已按「开始者写完」补完，见文末 [边界] 标记）。

% [原文 p.85]
我们由此得到 $(A,G)$ 理论的一个 Class S 描述：它来自把 $G$ 型的 6d $\cN=(2,0)$ SCFT 紧化在一个带非正则穿刺点的球面上（细节以及更多这类 IIB 与 Class S 对应的例子见 [47]）。不过，由其他类型 $(D,D)$、$(D,E)$ 与 $(E,E)$ 的三重复形奇点工程化出来的 SCFT，则都还没有任何已知的 Class S 描述。\footnote{反过来，任何 Class S 构造在 IIB 型中都有一个对偶的几何描述，即某个非紧 Calabi--Yau 三重复形。然而一般而言，这个三重复形中的奇点不会是孤立的。相关的评论见下一节以及 (5.32) 处的脚注。近期的讨论另见 [278,279]。}

% [Tutor Note: 原文此处说「下一节与脚注 64」，指的是下一节挂在 (5.32) 上的那条脚注（讨论 $(G,G')$ 型的孤立奇点在 T-对偶下的像）。因本译本脚注由 LaTeX 自动编号、与原文编号不必相同，这里改按内容指位，不写死「脚注 64」。]

在结束本小节之前，我们想强调：有了这样的等价关系，Class S 构造的全部数据都可以从对偶的 IIB 几何工程中的局部 ALE 纤维化里读取出来，其中也包括整体结构（global structures）——后者近来已在 [278,279] 等例子上被研究过。

\SUB{5.4}{进一步推广}

\SUBSUB{5.4.1}{一般的奇点}

在前面的各节中，我们主要考虑由某些孤立超曲面奇点所描述的非紧 Calabi--Yau 三重复形，并看到 IIB 型中的这类几何工程能给出许多有趣的 4d $\cN=2$ SCFT。满足定义 4d SCFT 所需条件的超曲面奇点，已在 [48,280] 中被完整列出（另见 [52,60]）。顺理成章的后续问题便是：能否越出 IHS 去工程化更多的 4d $\cN=2$ SCFT？

% [Tutor Note: IHS = isolated hypersurface singularity（孤立超曲面奇点），即 $\C^4$ 中由单个方程 $f(x_1,\ldots,x_4)=0$ 定义、且只有孤立奇点的超曲面，本综述前文已引入。它的「权重判据」是 $\sum_i q_i>1$（$q_i$ 为坐标在 $\C^*$ 作用下的权、按奇点齐次多项式的次数归一），而 4d $\cN=2$ 要求该 $\C^*$ 作用对应 $\mathrm{U}(1)_r$，故 $\sum_i q_i>1$ 正是 $r$ 荷为正的体现。下面第 4 条对 ICIS 的 $\sum_i w_i>\sum_i d_i$ 是同一判据的推广。]

事实上，人们也可以考虑更一般的孤立典范奇点，例如孤立完全交奇点（isolated complete intersection singularity，ICIS），它由若干个多项式定义为
\begin{equation}
W := F_1(x_1,x_2,\ldots,x_{n+3}) = F_2(x_1,x_2,\ldots,x_{n+3}) = \ldots = F_n(x_1,x_2,\ldots,x_{n+3}) = 0 \subset \C^{n+3}\,.
\tag{5.28}
\end{equation}

% [Tutor Note: (5.28) 的维数账：环境空间 $\C^{n+3}$ 复维 $n+3$，$n$ 个独立方程（Jacobi 矩阵秩为 $n$）把维数降 $n$，剩下的正是复 3 维（三重复形）；$n=1$ 时回到 IHS 的单个方程。]
要工程化一个 4d SCFT，奇点的结构须满足某些（充分）条件，它们与我们在 IHS 情形下讨论过的条件类似：
\begin{enumerate}
  \item 奇点必须是孤立的。也就是说，方程组 $F_1=F_2=\ldots=F_n=0$ 与 $dF_a=0$（$a=1,2,\ldots,n$）有唯一解。
  \item $W$ 是一个完全交，即 Jacobi 矩阵 $\frac{\partial F_a}{\partial x_i}$（$a=1,\ldots,n$，$\alpha=1,\ldots,n+3$）在除奇点以外处处秩为 $n$。
% [原文 p.86]
% [Tutor Note: 原文第 2 条把 Jacobi 矩阵写成 $\partial F_a/\partial x_i$，紧接着却把第二个指标记作 $\alpha=1,\ldots,n+3$（已用 600 dpi 裁图逐字符复核，确为 $x_i$ 与 $\alpha$）。两者不自洽：正确的写法应是 $\partial F_a/\partial x_\alpha$，$a=1,\ldots,n$、$\alpha=1,\ldots,n+3$，即一个 $n\times(n+3)$ 矩阵。译文照原文抄录，此处按其应有之义理解。]
  \item $W$ 中的每个多项式 $F_i$ 都是齐次的，次数为 $d_i$；而各坐标 $x_i$ 带有权 $w_i$。
  \item 这些权 $w_i$ 须满足条件 $\sum_i w_i > \sum_i d_i$，这与前面为 IHS 引入的 $\sum_i q_i>1$ 是同一个精神。
\end{enumerate}

最近 [50] 证明：上述条件要求 $n=2$，即 ICIS 最多由两个多项式 $(F_1,F_2)$ 定义，于是总共得到 303 类奇点。ICIS 的构造在许多方面与 IHS 相似。例如，ICIS 的 mini-versal 形变定义为
\begin{equation}
F = F + \sum_{\alpha=1}^{\mu} \lambda_\alpha \phi_\alpha\,.
\tag{5.29}
\end{equation}

% [Tutor Note: (5.29) 左右两端在原文中同为数学斜体的 $F$（用 PyMuPDF 取字体：两端 span 的 font 均为 CMMI12，即 \mathnormal 的 $F$，不是 \mathcal）。按下一句「$F:=(F_1,F_2)$ 是 $2\times1$ 列向量」的用法，左端应是形变后的对象、右端是未形变的 $F$，故左端本应另有记号（如 $\mathcal{F}$）。译文照原文抄录，读者按「左端为形变后」理解即可。]
这里 $F := (F_1,F_2)$ 是一个 $2\times 1$ 列向量，而 $\phi_\alpha$ 取作只有一项非零的 $2\times 1$ 列向量，它们给出 Milnor 环的一组单项式基
\begin{equation}
J(F_i) = \C^5/(dF)\,.
\tag{5.30}
\end{equation}

% [Tutor Note: (5.30) 中 $\C^5$ 的 5 来自 $n=2$ 时的 $n+3=5$，即 ICIS 的环境空间 $\C^5$；$(dF)$ 表示由 $F_1,F_2$ 的所有偏导生成的理想，于是 $J(F_i)$ 是 $\C^5$ 对该理想的商环（Milnor 环）。其复维数即 (5.29) 中的 $\mu$，也就是形变参数的个数。]
此处三重复形 ICIS 上的典范全纯 3-形式定义为
\begin{equation}
\Omega = \frac{dx_1 \wedge dx_2 \wedge \ldots \wedge dx_5}{dF_1 \wedge dF_2}\,,
\tag{5.31}
\end{equation}
它在所得的 4d SCFT 中约化为 4d $\cN=2$ 的 SW 微分。类似地，人们可以从 ICIS 的几何中读出各种物理量，例如诸参数 $\lambda$ 的标度维数、中心荷 $(a,c)$，正如从前几小节详述的 IHS 情形中读出的一样。

% [Tutor Note: (5.31) 的分子是 $\C^5$ 上的全纯 5-形式（$n+3=5$），分母是两个方程的全纯部分 $dF_1\wedge dF_2$，商在 ICIS 上良定义，给出复 3 维的 $(3,0)$ 形式；这正是「典范 3-形式」。它之所以能约化为 SW 微分，是因为 $\Omega$ 在 3-链上的周期正是 Coulomb 分支有效理论的 $(a_D,a)$——即 SW 微分的几何来源。]
此外，[48] 猜测：能生成 4d $\cN=2$ SCFT 的最一般的孤立奇点是\textbf{有理分次 Gorenstein}奇点。\footnote{有理分次 Gorenstein 奇点也称为指标一（index one）典范奇点。}这里的分次条件意味着奇点带有一个 $\C^*$ 作用，这正是 4d $\cN=2$ SCFT 的 $\mathrm{U}(1)_r$ 对称性所要求的；Gorenstein 条件意味着存在一个可区分的顶形式 $\Omega$，它约化为 4d $\cN=2$ 的 SW 微分；而有理性条件则保证顶形式 $\Omega$ 在 $\mathrm{U}(1)_r$ 对称性下具有正电荷。

% [Tutor Note: 三个条件各有分工，别混：分次（graded）$\Rightarrow$ 有 $\C^*$ 作用（对应 $\mathrm{U}(1)_r$，使算子有确定的 $r$ 荷）；Gorenstein $\Rightarrow$ 存在处处非零的顶形式 $\Omega$（对应存在 SW 微分，即 Coulomb 分支有效理论）；有理（rational）$\Rightarrow$ 该顶形式的 $r$ 荷为正（否则 $\Omega$ 作为「好」的算子不会出现在幺正谱里）。三者合起来正是 Fano/典范奇点分类里的 index-one 典范奇点。]

人们甚至可以越出孤立奇点，在几何工程中考虑非孤立奇点。此时可以出现复一维的奇异轨迹。事实上，正如我们在第 4 节中遇到的，Class S 构造可以被看成是
% [原文 p.87]
这一类，因为它们部分地由 ALE 空间的 ADE 奇点（而不是 Calabi--Yau 三重复形）来分类，从而把维数降低了一。于是一个一般的 Class S 构造可以改写为如下超曲面奇点所定义的形式
\begin{equation}
F(x_1,x_2,x_3,x_4) = F_{ADE}(w_1(x_i), w_2(x_i), w_3(x_i))\,,
\tag{5.32}
\end{equation}
其中 $F_{ADE}(w_1(z_i),w_2(z_i),w_3(z_i))$ 指的是以三个变量 $w_i$ 写出的一个 ADE 奇点。于是轨迹 $\{w_i=0\}$ 在由 $x_i$ 参数化的环境空间中给出一个曲线奇点。\footnote{好奇的读者不妨把它与 5.3 节讨论的 T-对偶作一比较：该对偶把某些 IHS 与 Class S 构造联系起来。事实上可以看到，对以 (5.32) 之形式写出的 $(G,G')$ 型孤立奇点，其 $\{w_i=0\}$ 轨迹是一个点而不是曲线。换个说法：经 T-对偶映射到 IIB 一侧时，一般的 Class S 构造在 IIB 的超曲面三重复形中给出的是曲线奇点，而 Class S 构造的一个特殊子集会给出孤立奇点（例如那些工程化 $(A,G)$ 理论的构造）。}不过，这类非孤立奇点的形变理论更为复杂，详细的讨论超出了本综述的范围。

% [Tutor Note: (5.32) 与紧随其后的句子在原文里不自洽：公式中三个函数写作 $w_1(x_i),w_2(x_i),w_3(x_i)$，紧接着的句子却写作 $w_1(z_i),w_2(z_i),w_3(z_i)$（已用 600 dpi 裁图复核，原文确为 $z_i$，不是 $x_i$ 的误认）。从下一句「由 $x_i$ 参数化的环境空间」可知 $z_i$ 应是 $x_i$。译文照原文抄录，读者按 $x_i$ 理解。]

\SUBSUB{5.4.2}{F 理论构造与 S-折}

既然 F 理论\footnote{关于 F 理论，我们推荐 TASI 讲义综述 [281]。}是 IIB 型弦论的非微扰完备化，并且它为量子场论提供了强有力的几何化，人们自然想知道：如何把 4d $\cN=2$ SCFT 直接嵌入 F 理论构造——这既可能有助于理解已知理论的精细特征，也可能产生此前未知的 SCFT。本节中我们想简要浏览这些构造，并看看 F 理论能为此添上哪些新的 4d 理论。

在 F 理论构造中，4d $\cN=2$ SCFT 通常被实现为一叠 D3 膜的世界体积理论，而这些 D3 膜在一叠轴子--伸缩子 $\tau$ 为常数的 $(p,q)$ 七维膜中作探针。这里 $(p,q)$ 在 IIB 型的 $\SL(2,\Z)$ 对偶下按二重态变换，而这些七维膜是熟悉的 D7 膜的强耦合表亲——后者对应于 $(1,0)$ 型。

这样一叠七维膜导致一个局部椭圆 K3 面\footnote{即 10d 背景几何局部为 $\R^{1,7}\times \mathbb{P}^1$，其中 $\mathbb{P}^1$ 是 K3 面的底。}，其 Kodaira 奇点类型以及规范群 $F$ 都是固定的；我们把它们连同对应的 $\tau$ 值一起列在表 6 中 [174,176]。沿 $\R^{1,3}\subset\R^{1,7}$ 延伸的那叠 D3 膜，可以看成七维膜上 8d SYM 理论的非阿贝尔瞬子——这是七维膜世界体积上的 Wess--Zumino 耦合 $\int \tr(F\wedge F)\wedge C_4$ 的结果。在这样的设置下，D3 膜堆的 4d 世界体积上承载一个 $\cN=2$ 规范理论，而轴子--伸缩子 $\tau$（即 K3 椭圆纤维的复结构）可以等同于 4d 规范耦合的复化。
% [边界] 本句在原文中跨 p.88→p.89，已在此写完整（「开始者写完」）；下一文件 tr08b 从 p.89 的第二句（「Since the seven-brane is non-compact …」）续写。

% [Tutor Note: 为何 $\tau$ 能当规范耦合：IIB 型中 $\tau=C_0+i e^{-\phi}$ 既是轴子--伸缩子，又在 D7 膜世界体积上以 $\tau=\theta/(2\pi)+4\pi i/g^2_{\rm YM}$ 的方式给出规范耦合；K3 的椭圆纤维模数就是 $\tau$，故「椭圆纤维的复结构」与「复化规范耦合」是同一个量。表 6 中 $\tau$ 取固定值，正说明这些点上理论是共形的（$\tau$ 不随模跑动）。]
